%%%%%%%%%%%%%%%%%%%%%%%%%%%%%%%%%%%%%%%%%%%%%%%%%%%%%%%%%%%%%%%%%%%%%%
% Lớp 8 | Chương II. Hằng đẳng thức đáng nhớ và ứng dụng | Lập phương của một tổng. Lập phương của một hiệu
% Dạng 8D2?7-2: Viết biểu thức dưới dạng lập phương
%%%%%%%%%%%%%%%%%%%%%%%%%%%%%%%%%%%%%%%%%%%%%%%%%%%%%%%%%%%%%%%%%%%%%%


% nguồn: dan_tu_AI (nhập 16/11/2026)
\begin{bt}%[8D2H7-2]%[ma:5884a64a]
	Viết biểu thức $x^3+9x^2+27x+27$ dưới dạng lập phương của một tổng.
	\loigiai{
		Ta có:
		\begin{align*} 
			x^3+9x^2+27x+27 &= x^3 + 3 \cdot x^2 \cdot 3 + 3 \cdot x \cdot 3^2 + 3^3 \\
			&= (x+3)^3.
		\end{align*}
	}
\end{bt}

% nguồn: dan_tu_AI (nhập 16/11/2026)
\begin{bt}%[8D2H7-2]%[ma:37ed1208]
	Viết biểu thức $8x^3-36x^2+54x-27$ dưới dạng lập phương của một hiệu.
	\loigiai{
		Ta có:
		\begin{align*} 
			8x^3-36x^2+54x-27 &= (2x)^3 - 3 \cdot (2x)^2 \cdot 3 + 3 \cdot (2x) \cdot 3^2 - 3^3 \\
			&= (2x-3)^3.
		\end{align*}
	}
\end{bt}

% nguồn: dan_tu_AI (nhập 16/11/2026)
\begin{bt}%[8D2H7-2]%[ma:05984d47]
	Viết biểu thức $x^3-6x^2y+12xy^2-8y^3$ dưới dạng lập phương của một hiệu.
	\loigiai{
		Ta có:
		\begin{align*} 
			x^3-6x^2y+12xy^2-8y^3 &= x^3 - 3 \cdot x^2 \cdot (2y) + 3 \cdot x \cdot (2y)^2 - (2y)^3 \\
			&= (x-2y)^3.
		\end{align*}
	}
\end{bt}

% nguồn: dan_tu_AI (nhập 16/11/2026)
\begin{bt}%[8D2H7-2]%[ma:44878500]
	Viết biểu thức $27x^3+27x^2y+9xy^2+y^3$ dưới dạng lập phương của một tổng.
	\loigiai{
		Ta có:
		\begin{align*} 
			27x^3+27x^2y+9xy^2+y^3 &= (3x)^3 + 3 \cdot (3x)^2 \cdot y + 3 \cdot (3x) \cdot y^2 + y^3 \\
			&= (3x+y)^3.
		\end{align*}
	}
\end{bt}

% nguồn: dan_tu_AI (nhập 16/11/2026)
\begin{bt}%[8D2H7-2]%[ma:b956e810]
	Viết biểu thức $\dfrac{1}{8}x^3 - \dfrac{3}{4}x^2 + \dfrac{3}{2}x - 1$ dưới dạng lập phương của một hiệu.
	\loigiai{
		Ta có:
		\begin{align*} 
			\dfrac{1}{8}x^3 - \dfrac{3}{4}x^2 + \dfrac{3}{2}x - 1 &= \left(\dfrac{1}{2}x\right)^3 - 3 \cdot \left(\dfrac{1}{2}x\right)^2 \cdot 1 + 3 \cdot \left(\dfrac{1}{2}x\right) \cdot 1^2 - 1^3 \\
			&= \left(\dfrac{1}{2}x - 1\right)^3.
		\end{align*}
	}
\end{bt}

% nguồn: dan_tu_AI (nhập 16/11/2026)
\begin{bt}%[8D2V7-2]%[ma:3260f813]
	Thể tích của một khối bê tông hình lập phương được cho bởi biểu thức $V = 64x^3 - 144x^2 + 108x - 27$ ($\text{dm}^3$). 
	\begin{enumerate}
		\item Tìm biểu thức tính độ dài cạnh khối bê tông đó.
		\item Tính độ dài cạnh của khối bê tông khi $x=2$.
	\end{enumerate}
	\loigiai{
		\begin{enumerate}
			\item Ta viết lại biểu thức tính thể tích $V$ dưới dạng lập phương của một hiệu:
			\begin{align*} 
				V &= 64x^3 - 144x^2 + 108x - 27 \\
				&= (4x)^3 - 3 \cdot (4x)^2 \cdot 3 + 3 \cdot (4x) \cdot 3^2 - 3^3 \\
				&= (4x-3)^3.
			\end{align*}
			Vì khối bê tông có dạng hình lập phương nên thể tích bằng lập phương độ dài cạnh. Do đó, biểu thức tính độ dài cạnh của khối bê tông là: $4x-3$ ($\text{dm}$).
			\item Khi $x=2$, độ dài cạnh của khối bê tông là:
			\[ 4 \cdot 2 - 3 = 8 - 3 = 5 \text{ (dm)}. \]
		\end{enumerate}
	}
\end{bt}
